\chapter{STrans 智慧交通视觉感知与管理系统设计方案 V2.0}\label{design-strans-ux667aux6167ux4ea4ux901aux89c6ux89c9ux611fux77e5ux4e0eux7ba1ux7406ux7cfbux7edfux8bbeux8ba1ux65b9ux6848-v2.0}

\section{文档信息}\label{design-ux6587ux6863ux4fe1ux606f}

\begin{longtable}[]{@{}
  >{\raggedright\arraybackslash}p{(\linewidth - 2\tabcolsep) * \real{0.5000}}
  >{\raggedright\arraybackslash}p{(\linewidth - 2\tabcolsep) * \real{0.5000}}@{}}
\toprule\noalign{}
\begin{minipage}[b]{\linewidth}\raggedright
项目
\end{minipage} & \begin{minipage}[b]{\linewidth}\raggedright
内容
\end{minipage} \\
\midrule\noalign{}
\endhead
\bottomrule\noalign{}
\endlastfoot
项目名称 & 基于多源视觉感知的智慧交通沙盘监测与分析系统 \\
软件名称 & STrans（Smart Transportation Sensing System） \\
文档名称 & 系统设计方案 \\
文档版本 & V2.0 \\
编制日期 & 2026-07-14 \\
项目阶段 & 结题交付 \\
事实基线 & \texttt{L\allowbreak{}I\allowbreak{}N\allowbreak{}G} 分支 \texttt{b0f3be2} 及当前结题工作区 \\
适用对象 & 开发、测试、部署维护人员及结题评审人员 \\
关联材料 & \texttt{09-\allowbreak{}系统设计说明书.\allowbreak{}md}、\texttt{13-\allowbreak{}软件工程图谱与\allowbreak{}P\allowbreak{}P\allowbreak{}T素材.\allowbreak{}md}、\texttt{14-\allowbreak{}初始要求追踪与答辩问题链.\allowbreak{}md}、\texttt{15-\allowbreak{}软件使用说明书.\allowbreak{}md}、\texttt{16-\allowbreak{}项目总结报告.\allowbreak{}md} \\
\end{longtable}

本文档描述当前交付系统的结构、关键机制、接口、数据、安全、部署和质量设计。文中``已实现''仅指能够由当前代码、测试或运行材料支撑的能力；``可选''\,``建议''\,``后续''不属于当前已完成主链路。

\subsection{版本变更记录}\label{design-ux7248ux672cux53d8ux66f4ux8bb0ux5f55}

\begin{longtable}[]{@{}
  >{\raggedright\arraybackslash}p{(\linewidth - 4\tabcolsep) * \real{0.3333}}
  >{\raggedright\arraybackslash}p{(\linewidth - 4\tabcolsep) * \real{0.3333}}
  >{\raggedright\arraybackslash}p{(\linewidth - 4\tabcolsep) * \real{0.3333}}@{}}
\toprule\noalign{}
\begin{minipage}[b]{\linewidth}\raggedright
版本
\end{minipage} & \begin{minipage}[b]{\linewidth}\raggedright
阶段
\end{minipage} & \begin{minipage}[b]{\linewidth}\raggedright
主要内容
\end{minipage} \\
\midrule\noalign{}
\endhead
\bottomrule\noalign{}
\endlastfoot
V1.0 & 初始设计 & 建立视频接入、后端服务、Web 展示和基础识别方案 \\
V2.0 & 结题设计 & 依据实际实现重构逻辑与部署架构；补充任务隔离、统一数据契约、道路建模、资源调度、数据闭环、安全、测试证据及限制；删除不再采用的输入方案；区分已验证能力与后续演进 \\
\end{longtable}

\subsection{设计范围}\label{design-ux8bbeux8ba1ux8303ux56f4}

V2.0 的交付范围是``多源视频接入 + 本地视觉分析 + 二维道路空间分析 + Web 管理与证据闭环''。系统输出白名单放行或拦截的软件建议，不控制物理执行机构；当前道路异常按``候选检测 + 人工复核''使用；独立碰撞识别、信号灯识别、完整跨线流量和三维数字孪生不纳入已完成范围。

\begin{center}\rule{0.5\linewidth}{0.5pt}\end{center}

\section{设计目标与原则}\label{design-ux8bbeux8ba1ux76eeux6807ux4e0eux539fux5219}

\subsection{设计目标}\label{design-ux8bbeux8ba1ux76eeux6807}

\begin{enumerate}
\def\labelenumi{\arabic{enumi}.}
\tightlist
\item
  对 RTSP/MJPEG、手机或 USB 摄像头、本地图片和录像进行统一接入与管理；
\item
  形成车辆检测、跟踪、车牌识别、白名单判断和交通状态分析主链路；
\item
  形成道路异物、道路行人和可选道路破损的候选检测链路；
\item
  将画面坐标映射为车道、路口、禁停区域和拥堵热力等道路语义；
\item
  建立认证授权、配置管理、历史查询、告警处置、证据留存和审计闭环；
\item
  根据计算资源和任务类型自适应调整推理参数，在 CUDA 不可用时允许 CPU 降级；
\item
  使系统结构、算法流程、测试结论和真实素材均可追踪、可复查。
\end{enumerate}

\subsection{设计原则}\label{design-ux8bbeux8ba1ux539fux5219}

\begin{itemize}
\tightlist
\item
  \textbf{纵向闭环优先}：优先保证``输入---分析---展示---处置---归档''端到端可运行；
\item
  \textbf{分层解耦}：视频、算法、道路逻辑、业务数据、API 和界面通过明确契约协作；
\item
  \textbf{任务隔离}：车辆监控与道路异常使用独立任务模式，防止异常结果污染车辆统计；
\item
  \textbf{统一契约}：识别模块统一产生 \texttt{Analysis\allowbreak{}Result}，减少前端、存储和算法服务之间的结构分叉；
\item
  \textbf{本地优先、可选扩展}：核心识别在本地完成，远程算法、转发和智能报告不作为主链路前提；
\item
  \textbf{可解释融合}：模型、几何、时序和业务规则共同形成结果，并保留失败原因的定位空间；
\item
  \textbf{安全与最小权限}：管理功能采用角色约束，凭据和视频地址避免进入公开输出；
\item
  \textbf{证据边界清晰}：测试通过率、系统可运行性和算法精度分别表述，不用样例输出代替真值指标。
\end{itemize}

\begin{center}\rule{0.5\linewidth}{0.5pt}\end{center}

\section{总体逻辑架构}\label{design-ux603bux4f53ux903bux8f91ux67b6ux6784}

系统采用分层、模块化的单机边缘架构，并预留网络化输入和可选外部服务接口。

\begin{longtable}[]{@{}
  >{\raggedright\arraybackslash}p{(\linewidth - 4\tabcolsep) * \real{0.3333}}
  >{\raggedright\arraybackslash}p{(\linewidth - 4\tabcolsep) * \real{0.3333}}
  >{\raggedright\arraybackslash}p{(\linewidth - 4\tabcolsep) * \real{0.3333}}@{}}
\toprule\noalign{}
\begin{minipage}[b]{\linewidth}\raggedright
层次
\end{minipage} & \begin{minipage}[b]{\linewidth}\raggedright
主要组成
\end{minipage} & \begin{minipage}[b]{\linewidth}\raggedright
主要职责
\end{minipage} \\
\midrule\noalign{}
\endhead
\bottomrule\noalign{}
\endlastfoot
展示层 & React、Vite、Canvas、浏览器语音能力 & 登录、实时监控、热力图、功能中心、道路建模入口 \\
API 与流服务层 & FastAPI、REST、MJPEG & 认证、配置、查询、推理、流式画面和管理接口 \\
业务服务层 & CameraHub、AlgorithmClient、RoadLogicService、AnalysisStore、AuthStore & 视频编排、算法调用、道路分析、业务持久化、权限与审计 \\
视觉感知层 & LocalModelService、RoadAnomalyService、RoadMaskService & 车辆/车牌、异常候选和道路语义分割 \\
资源与调度层 & SystemMonitor、AdaptiveModelScheduler & 采集资源指标并选择推理档位 \\
数据持久层 & SQLite、证据文件、道路模型 JSON、模型权重 & 保存配置、结果、事件、证据、审计和算法资产 \\
输入与可选外部服务 & 网络视频、本地媒体、远程算法、智能报告服务 & 提供数据源或非核心增强能力 \\
\end{longtable}

核心调用路径如下：

\begin{Verbatim}[breaklines=true,breakanywhere=true,fontsize=\small,frame=single,framesep=2mm,rulecolor=\color{STransLine}]
视频源
  → CameraHub / VideoStreamService
  → AdaptiveModelScheduler
  → LocalModelService 或 RoadAnomalyService
  → RoadMaskService / RoadLogicService
  → AnalysisResult
  → FastAPI + React 展示
  → SQLite 历史、告警、证据与审计
\end{Verbatim}

\pandocbounded{\includegraphics[keepaspectratio,alt={STrans 系统组件架构}]{assets/component-architecture.png}}

\emph{图 3-1　STrans 系统组件架构图}

该图强调``视频接入---视觉感知---道路逻辑---业务存储---API---前端''的组件边界。算法模块只通过统一分析结果和明确服务接口向上提供能力，道路建模工具则位于配置生产侧，不进入逐帧推理调用，从而降低界面、算法和空间模型之间的直接耦合。

组件边界和依赖关系见：

\begin{itemize}
\tightlist
\item
  \href{diagrams/component-architecture.puml}{组件架构 PlantUML 图源}
\item
  \href{../../output/ppt-assets/diagrams/component-architecture.svg}{组件架构 SVG}
\item
  \href{diagrams/realtime-analysis-sequence.puml}{实时分析时序图源}
\item
  \href{../../output/ppt-assets/diagrams/realtime-analysis-sequence.svg}{实时分析时序 SVG}
\end{itemize}

\pandocbounded{\includegraphics[keepaspectratio,alt={STrans 实时分析时序}]{assets/realtime-analysis-sequence.png}}

\emph{图 3-2　实时视频分析主链路时序图}

该时序图从用户启动视频源开始，展示取帧、资源调度、车辆或异常分析、道路语义补充、结果持久化以及浏览器接收标注流的先后关系。它说明 MJPEG 画面和结构化分析结果由同一摄像头上下文关联，但车辆监控与道路异常在算法分支和输出语义上保持隔离。

\begin{center}\rule{0.5\linewidth}{0.5pt}\end{center}

\section{部署架构}\label{design-ux90e8ux7f72ux67b6ux6784}

\subsection{已验证部署形态}\label{design-ux5df2ux9a8cux8bc1ux90e8ux7f72ux5f62ux6001}

当前结题验证以一台 Windows 11 主机为核心：FastAPI、SQLite、视觉模型和资源监控运行在本地主机，React 前端由 Vite 开发服务器或构建后的静态文件提供，浏览器作为操作终端；NVIDIA GPU 承担正式识别推理，CPU 可执行轻量测试和降级推理。

\begin{longtable}[]{@{}
  >{\raggedright\arraybackslash}p{(\linewidth - 4\tabcolsep) * \real{0.3333}}
  >{\raggedright\arraybackslash}p{(\linewidth - 4\tabcolsep) * \real{0.3333}}
  >{\raggedright\arraybackslash}p{(\linewidth - 4\tabcolsep) * \real{0.3333}}@{}}
\toprule\noalign{}
\begin{minipage}[b]{\linewidth}\raggedright
节点
\end{minipage} & \begin{minipage}[b]{\linewidth}\raggedright
进程或资源
\end{minipage} & \begin{minipage}[b]{\linewidth}\raggedright
默认端口/位置
\end{minipage} \\
\midrule\noalign{}
\endhead
\bottomrule\noalign{}
\endlastfoot
浏览器终端 & React Web 界面、MJPEG 展示、Canvas 热力图 & \texttt{http:\allowbreak{}/\allowbreak{}/\allowbreak{}127.\allowbreak{}0.\allowbreak{}0.\allowbreak{}1:\allowbreak{}5173} \\
应用主机 & FastAPI、视频线程、识别与道路逻辑 & \texttt{127.\allowbreak{}0.\allowbreak{}0.\allowbreak{}1:\allowbreak{}8000} \\
GPU & YOLO 等 CUDA 推理 & 主机本地设备 \\
数据层 & SQLite、证据、模型权重、道路模型 & 主机本地文件系统 \\
道路建模桥接 & 可选 RTSP 单帧提取 & 仅监听 \texttt{127.\allowbreak{}0.\allowbreak{}0.\allowbreak{}1:\allowbreak{}8765} \\
\end{longtable}

\subsection{可选扩展形态}\label{design-ux53efux9009ux6269ux5c55ux5f62ux6001}

网络摄像头、手机视频或媒体转发服务可经局域网/公网向应用主机提供 RTSP/MJPEG；远程算法服务和智能报告服务也可通过接口接入。上述能力具备接口和部署结构，但当前尚无三台独立物理设备的长期协同验收记录，因此不视为已完成的分布式生产部署。

部署节点和通信关系见：

\begin{itemize}
\tightlist
\item
  \href{diagrams/deployment-architecture.puml}{部署架构 PlantUML 图源}
\item
  \href{../../output/ppt-assets/diagrams/deployment-architecture.svg}{部署架构 SVG}
\end{itemize}

\pandocbounded{\includegraphics[keepaspectratio,alt={STrans 系统部署架构}]{assets/deployment-architecture.png}}

\emph{图 4-1　STrans 已验证部署与可选外部连接图}

图中以 Windows 应用主机为当前验证中心，浏览器、FastAPI、SQLite、本地模型和 NVIDIA GPU 构成已运行链路；摄像头、媒体转发、远程算法和智能报告按外部连接表示。该边界用于避免把``接口可接入''扩大为``多节点生产环境已验收''。

\begin{center}\rule{0.5\linewidth}{0.5pt}\end{center}

\section{功能模块设计}\label{design-ux529fux80fdux6a21ux5757ux8bbeux8ba1}

\begin{longtable}[]{@{}
  >{\raggedright\arraybackslash}p{(\linewidth - 8\tabcolsep) * \real{0.2000}}
  >{\raggedright\arraybackslash}p{(\linewidth - 8\tabcolsep) * \real{0.2000}}
  >{\raggedright\arraybackslash}p{(\linewidth - 8\tabcolsep) * \real{0.2000}}
  >{\raggedright\arraybackslash}p{(\linewidth - 8\tabcolsep) * \real{0.2000}}
  >{\raggedright\arraybackslash}p{(\linewidth - 8\tabcolsep) * \real{0.2000}}@{}}
\toprule\noalign{}
\begin{minipage}[b]{\linewidth}\raggedright
编号
\end{minipage} & \begin{minipage}[b]{\linewidth}\raggedright
模块
\end{minipage} & \begin{minipage}[b]{\linewidth}\raggedright
输入
\end{minipage} & \begin{minipage}[b]{\linewidth}\raggedright
输出/职责
\end{minipage} & \begin{minipage}[b]{\linewidth}\raggedright
当前状态
\end{minipage} \\
\midrule\noalign{}
\endhead
\bottomrule\noalign{}
\endlastfoot
M01 & 视频接入 & 视频地址、文件、图片、设备索引 & 统一帧、最新 JPEG、状态与重连 & 已实现并有测试/真实流材料 \\
M02 & 车辆感知 & 视频帧、推理配置 & 车辆框、类别、置信度、轨迹 & 已实现；小目标和遮挡仍有限制 \\
M03 & 车牌与白名单 & 车辆区域、全帧、轨迹、名单 & 稳定车牌文本、\texttt{allow/\allowbreak{}deny} 软件建议 & 已实现软件闭环；无实体控制 \\
M04 & 道路异常 & 道路 ROI、帧序列、合法目标框 & 异物/行人/路损候选 & 已实现候选链路；需人工复核 \\
M05 & 道路分割 & 画面帧 & 道路二值掩膜 & 已实现；模型资产需预先准备 \\
M06 & 道路逻辑 & 检测、标定、道路模型 & 世界坐标、车道/路口、速度、拥堵、禁停 & 已实现核心逻辑；非参考速度为估计 \\
M07 & 数据闭环 & \texttt{Analysis\allowbreak{}Result}、事件和证据帧 & 历史、告警、证据、导出 & 已实现；事件合并仍需增强 \\
M08 & 认证授权 & 用户、验证码、会话 & 登录、角色权限、审计 & 已实现基础 RBAC \\
M09 & 摄像头管理 & 视频源配置 & CRUD、测试、启停和状态 & 已实现视频源管理 \\
M10 & 智能报告 & 历史与事件摘要、外部配置 & 报告内容和归档 & 可选；依赖外部服务 \\
M11 & 监控与调度 & CPU、内存、GPU、显存、耗时、任务 & 档位和推理参数 & 已实现策略与页面；长稳不足 \\
M12 & Web 展示 & REST、MJPEG、轮询数据 & 实时画面、状态、管理页面 & 已实现并完成浏览器遍历 \\
M13 & 道路建模工具 & 节点、车道、建筑、摄像头、标定点 & \texttt{road\_\allowbreak{}logic\_\allowbreak{}modeler.\allowbreak{}v1} JSON & 已集成；发布仍需人工审核 \\
\end{longtable}

用户和管理员的系统边界见：

\begin{itemize}
\tightlist
\item
  \href{diagrams/use-case-overview.puml}{精简用例图源} / \href{../../output/ppt-assets/diagrams/use-case-overview.svg}{SVG}
\item
  \href{diagrams/use-case.puml}{完整用例图源} / \href{../../output/ppt-assets/diagrams/use-case.svg}{SVG}
\end{itemize}

\begin{center}\rule{0.5\linewidth}{0.5pt}\end{center}

\section{数据设计}\label{design-ux6570ux636eux8bbeux8ba1}

\subsection{统一分析契约}\label{design-ux7edfux4e00ux5206ux6790ux5951ux7ea6}

\texttt{Analysis\allowbreak{}Result} 是实时展示、道路逻辑、远程算法适配和持久化之间的核心契约，主要字段包括：

\begin{longtable}[]{@{}ll@{}}
\toprule\noalign{}
字段组 & 主要内容 \\
\midrule\noalign{}
\endhead
\bottomrule\noalign{}
\endlastfoot
帧上下文 & \texttt{frame\_\allowbreak{}id}、\texttt{timestamp}、\texttt{camera\_\allowbreak{}id}、源尺寸 \\
模型上下文 & \texttt{model\_\allowbreak{}id}、\texttt{inference\_\allowbreak{}ms}、错误信息 \\
检测结果 & 边界框、类别、置信度、\texttt{track\_\allowbreak{}id}、车牌、白名单状态、速度 \\
交通状态 & 活跃车辆数、密度、拥堵等级、热力数据 \\
事件 & \texttt{event\_\allowbreak{}id}、类型、等级、描述、位置和时间 \\
扩展元数据 & 任务模式、调度档位、道路语义等可扩展内容 \\
\end{longtable}

车辆监控和道路异常虽然共用该结构，但应保持语义隔离：异常模式不将异常候选计入车辆交通统计，车辆模式不混入道路异常结果。

\subsection{持久化模型}\label{design-ux6301ux4e45ux5316ux6a21ux578b}

\begin{longtable}[]{@{}
  >{\raggedright\arraybackslash}p{(\linewidth - 4\tabcolsep) * \real{0.3333}}
  >{\raggedright\arraybackslash}p{(\linewidth - 4\tabcolsep) * \real{0.3333}}
  >{\raggedright\arraybackslash}p{(\linewidth - 4\tabcolsep) * \real{0.3333}}@{}}
\toprule\noalign{}
\begin{minipage}[b]{\linewidth}\raggedright
数据类别
\end{minipage} & \begin{minipage}[b]{\linewidth}\raggedright
主要表/产物
\end{minipage} & \begin{minipage}[b]{\linewidth}\raggedright
关联方式
\end{minipage} \\
\midrule\noalign{}
\endhead
\bottomrule\noalign{}
\endlastfoot
分析历史 & \texttt{analysis\_\allowbreak{}records} & 摄像头、时间、检测/事件 JSON、耗时 \\
告警事件 & \texttt{alert\_\allowbreak{}incidents} & 以 \texttt{event\_\allowbreak{}id} 标识和处置 \\
证据帧 & \texttt{evidence\_\allowbreak{}frames}、证据文件 & 通过 \texttt{event\_\allowbreak{}id} 关联，记录 SHA-256 摘要 \\
配置 & \texttt{camera\_\allowbreak{}sources}、\texttt{vehicle\_\allowbreak{}whitelist}、\texttt{intelligence\_\allowbreak{}config} & 由管理接口维护 \\
用户安全 & \texttt{system\_\allowbreak{}users}、\texttt{system\_\allowbreak{}sessions}、\texttt{audit\_\allowbreak{}logs} & 用户、会话与操作关联 \\
道路模型 & 版本化 JSON 文件 & 摄像头标定和道路实体引用 \\
\end{longtable}

SQLite 符合当前单机、低运维和便于备份的交付条件。它不适合未经评估的多实例并发写入；若后续转为多节点高并发服务，应重新评估事务、迁移和数据库选型。

数据闭环见 \href{diagrams/data-closed-loop-sequence.puml}{时序图源} 和 \href{../../output/ppt-assets/diagrams/data-closed-loop-sequence.svg}{SVG}。

\pandocbounded{\includegraphics[keepaspectratio,alt={STrans 分析数据闭环}]{assets/data-closed-loop-sequence.png}}

\emph{图 6-1　分析记录、告警、证据与处置闭环时序图}

该图以 \texttt{event\_\allowbreak{}id} 为业务关联主线，说明分析结果如何进入历史记录，候选事件如何关联原始帧、标注帧和结构化数据，以及用户如何完成查询、处置和证据导出。SHA-256 摘要用于核对导出证据完整性；事件业务级去重和生命周期合并仍是后续增强项。

\begin{center}\rule{0.5\linewidth}{0.5pt}\end{center}

\section{API 与交互设计}\label{design-api-ux4e0eux4ea4ux4e92ux8bbeux8ba1}

\subsection{接口分组}\label{design-ux63a5ux53e3ux5206ux7ec4}

\begin{longtable}[]{@{}
  >{\raggedright\arraybackslash}p{(\linewidth - 4\tabcolsep) * \real{0.3333}}
  >{\raggedright\arraybackslash}p{(\linewidth - 4\tabcolsep) * \real{0.3333}}
  >{\raggedright\arraybackslash}p{(\linewidth - 4\tabcolsep) * \real{0.3333}}@{}}
\toprule\noalign{}
\begin{minipage}[b]{\linewidth}\raggedright
分组
\end{minipage} & \begin{minipage}[b]{\linewidth}\raggedright
代表接口
\end{minipage} & \begin{minipage}[b]{\linewidth}\raggedright
设计目的
\end{minipage} \\
\midrule\noalign{}
\endhead
\bottomrule\noalign{}
\endlastfoot
健康与认证 & \texttt{/\allowbreak{}api/\allowbreak{}health}、\texttt{/\allowbreak{}api/\allowbreak{}auth/\allowbreak{}captcha}、\texttt{/\allowbreak{}api/\allowbreak{}auth/\allowbreak{}login}、\texttt{/\allowbreak{}api/\allowbreak{}auth/\allowbreak{}me} & 健康检查、验证码、登录和会话 \\
摄像头 & \texttt{/\allowbreak{}api/\allowbreak{}cameras}、\texttt{/\allowbreak{}\{id\}/\allowbreak{}test}、\texttt{/\allowbreak{}\{id\}/\allowbreak{}start}、\texttt{/\allowbreak{}\{id\}/\allowbreak{}stop} & 视频源管理和生命周期 \\
视频流 & \texttt{/\allowbreak{}\{id\}/\allowbreak{}mjpeg}、\texttt{/\allowbreak{}\{id\}/\allowbreak{}model-\allowbreak{}mjpeg} & 原始或标注 MJPEG 输出 \\
算法 & \texttt{/\allowbreak{}api/\allowbreak{}algorithm/\allowbreak{}config}、\texttt{/\allowbreak{}api/\allowbreak{}algorithm/\allowbreak{}infer/\allowbreak{}\{camera\_\allowbreak{}id\}} & 参数配置和单帧推理 \\
实时分析 & \texttt{/\allowbreak{}api/\allowbreak{}analysis/\allowbreak{}latest}、\texttt{/\allowbreak{}api/\allowbreak{}analysis/\allowbreak{}events} & 最新结果和事件查询 \\
业务闭环 & \texttt{/\allowbreak{}api/\allowbreak{}history}、\texttt{/\allowbreak{}api/\allowbreak{}incidents}、证据与导出接口 & 查询、处置、导出和追溯 \\
道路与资源 & 道路热力、系统资源、调度配置接口 & 道路状态和运行策略展示 \\
\end{longtable}

\subsection{交互约定}\label{design-ux4ea4ux4e92ux7ea6ux5b9a}

\begin{itemize}
\tightlist
\item
  REST 接口使用 JSON；模型输出遵循统一分析契约；
\item
  视频采用 \texttt{multipart/\allowbreak{}x-\allowbreak{}mixed-\allowbreak{}replace} MJPEG，浏览器通过 \texttt{<img>} 持续接收；
\item
  分析结果、摄像头状态、资源和热力图采用约 1.5～2 Hz 的短轮询；
\item
  管理类写操作由角色依赖执行权限检查，并写入审计；
\item
  外部算法和报告服务失败不应中断本地视频与核心分析链路；
\item
  当前设计未引入 WebSocket，若后续并发和实时性要求提高，应先完成负载测量再决定迁移。
\end{itemize}

\begin{center}\rule{0.5\linewidth}{0.5pt}\end{center}

\section{安全设计}\label{design-ux5b89ux5168ux8bbeux8ba1}

\subsection{已实现措施}\label{design-ux5df2ux5b9eux73b0ux63aaux65bd}

\begin{itemize}
\tightlist
\item
  登录与注册使用限时图片验证码；
\item
  会话使用随机令牌，支持当前用户查询、修改密码和主动登出；
\item
  采用 admin/user 两级 RBAC 限制摄像头、白名单、模型、用户和审计等管理操作；
\item
  视频地址在日志和审计输出中脱敏；外部服务密钥读取时仅显示脱敏信息；
\item
  告警证据记录 SHA-256 摘要，用于导出后的完整性核验；
\item
  管理操作进入审计日志，记录操作人、动作、目标和时间；
\item
  道路建模桥接默认只监听本机回环地址。
\end{itemize}

\subsection{安全边界与改进}\label{design-ux5b89ux5168ux8fb9ux754cux4e0eux6539ux8fdb}

当前认证和密码存储方案适合课程交付环境，不应直接等同于互联网生产级安全。面向生产部署时，应采用专用慢哈希算法、HTTPS、令牌吊销与刷新策略、登录限流、更严格的 CORS/CSRF 策略、密钥托管、依赖漏洞扫描和备份加密。公开材料不得包含完整视频凭据、管理员密码、外部 API 密钥、个人信息或未脱敏日志。

\begin{center}\rule{0.5\linewidth}{0.5pt}\end{center}

\section{算法设计}\label{design-ux7b97ux6cd5ux8bbeux8ba1}

\subsection{车辆、跟踪、车牌与白名单}\label{design-ux8f66ux8f86ux8ddfux8e2aux8f66ux724cux4e0eux767dux540dux5355}

车辆链路采用``检测---跟踪---稳态增强---业务判断''四段式设计：

\begin{enumerate}
\def\labelenumi{\arabic{enumi}.}
\tightlist
\item
  对输入帧执行尺寸调整和道路 ROI 约束，YOLO 产生车辆候选；
\item
  ByteTrack 维护跨帧身份，并结合重叠抑制、视觉恢复、框平滑和短时预测改善低帧率与遮挡场景；
\item
  对车辆区域执行 HyperLPR3，周期性全帧 OCR 作为补充，使用预算轮转、文本纠正、轨迹绑定和多帧投票稳定结果；
\item
  将稳定车牌与白名单匹配，输出 \texttt{allow/\allowbreak{}deny} 软件建议，并关联历史、事件和证据；
\item
  道路逻辑根据轨迹、标定和道路模型计算速度估计、车道/路口归属、禁停和拥堵状态。
\end{enumerate}

流程见 \href{diagrams/vehicle-recognition-activity.puml}{车辆识别活动图源} 和 \href{../../output/ppt-assets/diagrams/vehicle-recognition-activity.svg}{SVG}。

\pandocbounded{\includegraphics[keepaspectratio,alt={车辆识别与业务判断活动流程}]{assets/vehicle-recognition-activity.png}}

\emph{图 9-1　车辆检测、跟踪、车牌稳定与业务判断活动图}

图中将单帧模型输出与跨帧稳态处理分开：YOLO 负责候选发现，ByteTrack 和视觉恢复保持目标身份，车牌 OCR 通过轨迹绑定与多帧投票形成稳定文本，最后再进入白名单和道路逻辑。这样可以把``识别到了什么''和``业务如何使用结果''分别验证，也能定位小目标、遮挡或 OCR 波动发生在哪个阶段。

\subsection{道路异常候选}\label{design-ux9053ux8defux5f02ux5e38ux5019ux9009}

由于沙盘异常真值规模不足，当前不使用单一端到端异常模型作为唯一判断，而采用可解释的多线索融合：

\begin{enumerate}
\def\labelenumi{\arabic{enumi}.}
\tightlist
\item
  在道路 ROI 内计算帧差、背景差分和稀疏光流；
\item
  判断相机运动和全局光照变化，避免将全画面变化解释为局部异物；
\item
  使用车辆/行人检测与近期合法轨迹框解释正常道路使用者；
\item
  综合亮度、Lab 色差、HSV 暖色、纹理、轮廓和面积规则过滤候选；
\item
  使用多帧 IoU 与确认计数提高稳定性，可选路损权重作为补充；
\item
  输出候选事件，交由界面和人工复核，不宣称低误报完备识别。
\end{enumerate}

流程见 \href{diagrams/road-anomaly-activity.puml}{道路异常活动图源} 和 \href{../../output/ppt-assets/diagrams/road-anomaly-activity.svg}{SVG}。

\pandocbounded{\includegraphics[keepaspectratio,alt={道路异常候选分析活动流程}]{assets/road-anomaly-activity.png}}

\emph{图 9-2　道路异常候选多线索融合活动图}

该图突出异常链路的``先解释、后判异''原则：先利用光流和全局变化判断相机或光照扰动，再以车辆、行人和近期轨迹解释合法道路使用者，最后对剩余局部变化执行外观过滤和多帧确认。输出因此定位为可追溯的候选事件，而不是未经人工复核的确定性结论。

\subsection{道路分割与空间映射}\label{design-ux9053ux8defux5206ux5272ux4e0eux7a7aux95f4ux6620ux5c04}

SegFormer-B0 产生道路掩膜，用于约束移动或未标定视角的画面热点。固定视角使用摄像头标定点和 RANSAC 单应矩阵，将车辆框底边中心投影到二维道路世界坐标，再判断道路有效性、车道、路口和禁停区域。该方法依赖摄像头姿态与道路模型版本一致；摄像头移动后必须重新标定。

\begin{center}\rule{0.5\linewidth}{0.5pt}\end{center}

\section{道路建模设计}\label{design-ux9053ux8defux5efaux6a21ux8bbeux8ba1}

道路建模工具负责生产空间配置，不参与逐帧神经网络推理。其输入包括世界尺寸、节点、节点组、平行车道、建筑物、摄像头以及画面点---道路实体标定关系，输出 \texttt{road\_\allowbreak{}logic\_\allowbreak{}modeler.\allowbreak{}v1} JSON 和派生逻辑。

当前发布流程为：

\begin{Verbatim}[breaklines=true,breakanywhere=true,fontsize=\small,frame=single,framesep=2mm,rulecolor=\color{STransLine}]
可视化编辑
  → 结构规范化与派生逻辑生成
  → 导出 road_logic_modeler.v1 JSON
  → 人工检查 schema、引用、方向、尺寸和标定点
  → 版本留存
  → RoadLogicService 加载使用
\end{Verbatim}

工具已随项目静态资源集成，并提供本机可选单帧桥接；但尚未实现完整的``草稿---审核---发布---回滚''服务端工作流，也不会自动覆盖生产配置。流程见 \href{diagrams/road-model-pipeline.puml}{道路建模图源} 和 \href{../../output/ppt-assets/diagrams/road-model-pipeline.svg}{SVG}。

\pandocbounded{\includegraphics[keepaspectratio,alt={道路建模到道路逻辑的配置流水线}]{assets/road-model-pipeline.png}}

\emph{图 10-1　道路建模、审核发布与运行消费流程图}

该图把编辑工具与运行时道路逻辑明确分离：工具生产可编辑模型和派生 JSON，人工审核负责控制配置质量，RoadLogicService 只消费已确认版本。单应映射把画面坐标转换为道路空间语义，因此摄像头标定、道路模型版本和运行配置必须成组留存。

\begin{center}\rule{0.5\linewidth}{0.5pt}\end{center}

\section{自适应调度设计}\label{design-ux81eaux9002ux5e94ux8c03ux5ea6ux8bbeux8ba1}

SystemMonitor 提供 CPU、内存、GPU、显存和推理耗时，AdaptiveModelScheduler 同时考虑任务模式、静态/实时来源和资源压力，选择以下档位：

\begin{longtable}[]{@{}lll@{}}
\toprule\noalign{}
档位 & 适用条件 & 主要策略 \\
\midrule\noalign{}
\endhead
\bottomrule\noalign{}
\endlastfoot
\texttt{quality} & 静态或质量优先、资源充足 & 较高输入尺寸与质量参数 \\
\texttt{balanced} & 常规实时分析 & 平衡尺寸、阈值和检测间隔 \\
\texttt{realtime} & 实时性优先或资源压力上升 & 降低尺寸/频率以控制延迟 \\
\texttt{protect} & 资源接近上限 & 立即降级，优先保持系统可用 \\
\texttt{anomaly} & 道路异常任务 & 保持适于小候选的任务参数 \\
\texttt{manual} & 管理员关闭自动调度 & 使用人工配置 \\
\end{longtable}

资源指标使用短期缓存；非保护档采用 8 秒滞回，减少瞬时波动引起的频繁切换；\texttt{protect} 不等待滞回。调度仅改变模型、输入尺寸、阈值和检测间隔，不改变事件业务定义。状态转换见 \href{diagrams/adaptive-scheduler-state.puml}{调度状态图源} 和 \href{../../output/ppt-assets/diagrams/adaptive-scheduler-state.svg}{SVG}。

\pandocbounded{\includegraphics[keepaspectratio,alt={自适应模型调度状态转换}]{assets/adaptive-scheduler-state.png}}

\emph{图 11-1　自适应模型调度状态图}

状态图将任务要求、实时性和资源压力统一映射为有限档位，避免在业务代码中散布参数切换规则。普通档位通过滞回抑制瞬时抖动，资源越界时则可立即进入 \texttt{protect}；\texttt{manual} 保留人工控制，但不会改变车辆、异常和告警的业务定义。

\begin{center}\rule{0.5\linewidth}{0.5pt}\end{center}

\section{可观测性与运行保障}\label{design-ux53efux89c2ux6d4bux6027ux4e0eux8fd0ux884cux4fddux969c}

\subsection{当前可观测信号}\label{design-ux5f53ux524dux53efux89c2ux6d4bux4fe1ux53f7}

\begin{longtable}[]{@{}lll@{}}
\toprule\noalign{}
信号 & 来源 & 用途 \\
\midrule\noalign{}
\endhead
\bottomrule\noalign{}
\endlastfoot
服务健康 & \texttt{/\allowbreak{}api/\allowbreak{}health}、模型健康状态 & 启动检查和模型可用性判断 \\
视频状态 & 摄像头状态、最新帧和失败重试 & 发现地址、解码和断流问题 \\
资源指标 & CPU、内存、GPU、显存 & 触发调度和定位资源瓶颈 \\
算法指标 & \texttt{inference\_\allowbreak{}ms}、模型、输入尺寸、档位 & 性能分析和环境复现 \\
业务记录 & 历史、事件、处置状态 & 追踪分析与业务结果 \\
审计记录 & 用户、动作、目标、时间 & 管理操作追责与复核 \\
证据摘要 & 原始/标注帧、JSON、SHA-256 & 结果核验和答辩材料归档 \\
\end{longtable}

\subsection{运维策略}\label{design-ux8fd0ux7ef4ux7b56ux7565}

\begin{itemize}
\tightlist
\item
  第一次正式识别前完成模型加载和 CUDA 预热；
\item
  每次演示记录 Python、PyTorch、CUDA、GPU、模型权重、阈值和道路模型版本；
\item
  定期备份 SQLite、道路模型、模型权重和证据目录；
\item
  视频源、道路模型或依赖变化后执行固定真实视频回归；
\item
  当前未形成集中式日志、指标时序数据库和告警平台，30～60 分钟长稳与断流恢复也尚未完成系统验收。
\end{itemize}

\begin{center}\rule{0.5\linewidth}{0.5pt}\end{center}

\section{测试与质量设计}\label{design-ux6d4bux8bd5ux4e0eux8d28ux91cfux8bbeux8ba1}

\subsection{已执行验证}\label{design-ux5df2ux6267ux884cux9a8cux8bc1}

\begin{longtable}[]{@{}
  >{\raggedright\arraybackslash}p{(\linewidth - 4\tabcolsep) * \real{0.3000}}
  >{\raggedleft\arraybackslash}p{(\linewidth - 4\tabcolsep) * \real{0.4000}}
  >{\raggedright\arraybackslash}p{(\linewidth - 4\tabcolsep) * \real{0.3000}}@{}}
\toprule\noalign{}
\begin{minipage}[b]{\linewidth}\raggedright
验证层次
\end{minipage} & \begin{minipage}[b]{\linewidth}\raggedleft
结果
\end{minipage} & \begin{minipage}[b]{\linewidth}\raggedright
能证明的范围
\end{minipage} \\
\midrule\noalign{}
\endhead
\bottomrule\noalign{}
\endlastfoot
后端 Python & 62/62 通过 & 认证、存储、白名单、摄像头、调度和道路逻辑等已有用例 \\
前端 JavaScript & 19/19 通过 & 前端已有组件与逻辑用例 \\
道路建模工具 & Python 21/21（含桥接 6 项）、Node 5/5 通过 & 建模、分析、桥接输入校验与本机约束 \\
浏览器系统遍历 & 17/17 检查点通过 & 双角色、主要页面和道路建模页面 \\
真实视频 CUDA & 13 段视频、65 帧全部成功读取和推理 & 数据链路、样例输出和当前环境性能 \\
\end{longtable}

上述合计 107 个自动化测试/脚本单元全部通过，但不表示需求、权限组合或算法准确率达到 100\%。当前真实视频没有完整人工框、车牌和事件真值，因而不得用检测输出数量代替 Precision、Recall、F1 或整牌准确率。

\subsection{质量门禁与后续测试}\label{design-ux8d28ux91cfux95e8ux7981ux4e0eux540eux7eedux6d4bux8bd5}

当前交付应至少执行 Python 测试、前端测试、生产构建、道路建模测试、PlantUML 校验和固定视频短回归。后续应补充：

\begin{enumerate}
\def\labelenumi{\arabic{enumi}.}
\tightlist
\item
  FastAPI 权限、参数、错误码和未认证访问矩阵；
\item
  视频断流、重连、损坏文件和多路并发测试；
\item
  事件冷却、合并和完整生命周期状态机测试；
\item
  禁停、白名单联动、道路热力和异常的专项真值集；
\item
  30～60 分钟长稳、资源压力和保护档切换测试；
\item
  车辆检测、车牌、异常和跟踪的分层精度指标。
\end{enumerate}

\begin{center}\rule{0.5\linewidth}{0.5pt}\end{center}

\section{部署与配置方案}\label{design-ux90e8ux7f72ux4e0eux914dux7f6eux65b9ux6848}

\subsection{环境基线}\label{design-ux73afux5883ux57faux7ebf}

结题验证环境为 Windows 11、Python 3.14.3、PyTorch 2.11.0+cu126、NVIDIA GeForce RTX 4050 Laptop GPU、Node.js 24.14.0 和 npm 11.12.1。正式识别使用 CUDA 环境；项目 CPU 虚拟环境可用于部分测试和基础检查。更换 Python、PyTorch、CUDA、GPU、模型或摄像头后，应重新预热、标定并执行回归。

\subsection{启动顺序}\label{design-ux542fux52a8ux987aux5e8f}

\begin{enumerate}
\def\labelenumi{\arabic{enumi}.}
\tightlist
\item
  检查模型权重、数据库、道路模型、证据目录和端口；
\item
  在 CUDA 可用的 Python 环境启动 FastAPI，访问 \texttt{/\allowbreak{}api/\allowbreak{}health}；
\item
  启动前端开发服务或部署生产构建产物；
\item
  使用管理员账号检查模型、调度、摄像头和道路模型配置；
\item
  测试视频源并启动一路固定样本完成预热；
\item
  核对标注画面、资源指标、历史和告警后进入正式演示；
\item
  演示结束后导出证据并备份数据库与配置。
\end{enumerate}

\subsection{配置治理}\label{design-ux914dux7f6eux6cbbux7406}

\begin{itemize}
\tightlist
\item
  管理员密码和外部服务密钥通过部署配置提供，不写入公开仓库；
\item
  录像路径、视频地址和道路模型应与目标主机环境一致；
\item
  道路模型、摄像头标定、模型权重和算法阈值作为同一运行基线记录；
\item
  SQLite 不由多个独立后端实例同时写入；
\item
  外部报告和远程算法未配置时，本地核心链路仍应可运行。
\end{itemize}

\begin{center}\rule{0.5\linewidth}{0.5pt}\end{center}

\section{关键设计决策}\label{design-ux5173ux952eux8bbeux8ba1ux51b3ux7b56}

\begin{longtable}[]{@{}
  >{\raggedright\arraybackslash}p{(\linewidth - 4\tabcolsep) * \real{0.3333}}
  >{\raggedright\arraybackslash}p{(\linewidth - 4\tabcolsep) * \real{0.3333}}
  >{\raggedright\arraybackslash}p{(\linewidth - 4\tabcolsep) * \real{0.3333}}@{}}
\toprule\noalign{}
\begin{minipage}[b]{\linewidth}\raggedright
决策
\end{minipage} & \begin{minipage}[b]{\linewidth}\raggedright
原因
\end{minipage} & \begin{minipage}[b]{\linewidth}\raggedright
影响
\end{minipage} \\
\midrule\noalign{}
\endhead
\bottomrule\noalign{}
\endlastfoot
车辆与异常任务隔离 & 两类任务目标、统计和调度不同 & 避免结果污染，但前端和测试需显式携带任务模式 \\
统一 \texttt{Analysis\allowbreak{}Result} & 降低算法、API、前端和存储耦合 & 新字段应保持向后兼容并有默认值 \\
SQLite 单机持久化 & 零配置、易备份，适合课程和沙盘部署 & 不适合未经评估的多实例高并发写入 \\
多线索异常候选 & 缺少大规模沙盘异常真值，要求可解释 & 可逐层定位误报，但需持续维护场景规则 \\
本地 CUDA 优先、CPU 降级 & 降低外部依赖并兼顾可运行性 & 设备和环境差异会改变时延与结果 \\
道路模型人工审核发布 & 避免错误标定直接影响运行分析 & 当前发布效率较低，后续需版本治理服务 \\
调度引入滞回 & 防止资源波动导致模型频繁切换 & 保护档立即生效，其他档位响应存在有意延迟 \\
\end{longtable}

算法工程演进依据见 \href{diagrams/algorithm-evolution-flow.puml}{算法演进图源} 和 \href{../../output/ppt-assets/diagrams/algorithm-evolution-flow.svg}{SVG}。

\begin{center}\rule{0.5\linewidth}{0.5pt}\end{center}

\section{已知限制与风险}\label{design-ux5df2ux77e5ux9650ux5236ux4e0eux98ceux9669}

\begin{longtable}[]{@{}
  >{\raggedright\arraybackslash}p{(\linewidth - 6\tabcolsep) * \real{0.2500}}
  >{\raggedright\arraybackslash}p{(\linewidth - 6\tabcolsep) * \real{0.2500}}
  >{\raggedright\arraybackslash}p{(\linewidth - 6\tabcolsep) * \real{0.2500}}
  >{\raggedright\arraybackslash}p{(\linewidth - 6\tabcolsep) * \real{0.2500}}@{}}
\toprule\noalign{}
\begin{minipage}[b]{\linewidth}\raggedright
风险/限制
\end{minipage} & \begin{minipage}[b]{\linewidth}\raggedright
当前影响
\end{minipage} & \begin{minipage}[b]{\linewidth}\raggedright
当前控制
\end{minipage} & \begin{minipage}[b]{\linewidth}\raggedright
后续方向
\end{minipage} \\
\midrule\noalign{}
\endhead
\bottomrule\noalign{}
\endlastfoot
无完整人工真值 & 不能给出严格算法精度 & 只报告样例、性能和失败案例 & 建立分层标注集和 Precision/Recall/F1 \\
道路箭头、标线和反光误报 & 产生道路异常候选 & 人工复核、保留失败样本 & 标线先验、长期背景和摄像头级阈值 \\
同类事件逐帧重复 & 告警和统计可能被淹没 & 演示期控制运行时间并人工处置 & 空间匹配、冷却、合并和生命周期 \\
极远景、遮挡、极小车牌 & 漏检或 OCR 波动 & ROI、跟踪、轮转、多帧稳定 & 专项数据与模型/镜头优化 \\
非参考摄像头速度为估计 & 不能视为法定测速 & 明确单位与口径 & 标定校验和真实速度对照 \\
摄像头或道路模型错配 & 车道、禁停和热力位置偏移 & 人工审核和版本留存 & 草稿、校验、发布与回滚 API \\
环境与设备差异 & 性能和部分数值结果变化 & 固定环境清单、预热和回归 & 容器化或锁定依赖矩阵 \\
长稳与断流覆盖不足 & 生产稳定性结论有限 & 已完成短时功能与浏览器验证 & 30～60 分钟长稳和故障注入 \\
外部服务不可用 & 报告或远程算法增强失败 & 本地核心链路独立运行 & 超时、重试、熔断和服务级监控 \\
基础认证非互联网生产级 & 直接公网部署风险较高 & 本机/受控网络、RBAC、脱敏和审计 & HTTPS、慢哈希、限流和密钥托管 \\
\end{longtable}

\begin{center}\rule{0.5\linewidth}{0.5pt}\end{center}

\section{设计验收准则}\label{design-ux8bbeux8ba1ux9a8cux6536ux51c6ux5219}

V2.0 设计是否落实，应以以下条件共同判断：

\begin{itemize}
\tightlist
\item
  主要模块、接口、数据和页面能够对应到当前实现；
\item
  车辆监控与道路异常模式的数据语义互不污染；
\item
  至少一路固定视频能闭合``启动---推理---展示---历史---告警''流程；
\item
  道路建模可以导出规范 JSON，并经人工审核后由道路逻辑消费；
\item
  CUDA 环境、模型、配置、道路模型和测试数据可以追溯；
\item
  管理操作遵循角色约束，敏感配置不出现在公开材料；
\item
  自动化、浏览器和真实视频证据分别说明覆盖范围，不扩大为全面精度或生产稳定性结论；
\item
  已知误报、事件重复、真值不足和长稳缺口在报告与答辩中保持同一口径。
\end{itemize}

\section{结论}\label{design-ux7ed3ux8bba}

STrans V2.0 将多源视频、视觉感知、道路空间模型、资源调度、Web 管理和数据证据组织为一个可运行的单机边缘系统。其核心设计价值是以统一数据契约连接模型、几何、时序和业务规则，并通过任务隔离、资源降级和证据闭环提高系统的可解释性、可管理性和可复查性。

当前设计已经覆盖结题主链路，但仍受真值不足、道路标线误报、事件合并缺失、设备差异和长稳测试不足等约束。V2.0 因此不以``算法完全准确''作为结论，而以``主链路可运行、设计边界清晰、证据可复查、后续改进可落地''作为交付口径。
